%mac

\begin{vidu}
Một vật có khối lượng $2\ kg$ đang nằm yên trên mặt sàn nằm ngang không ma sát. Một lực kéo không đổi $F = 6\ N$ bắt đầu tác dụng lên vật theo phương ngang. Sau khi chịu lực trong $3\ s$, vật đạt được tốc độ $v$ và đi được quãng đường $s$.
\begin{listEX}
\item Tính gia tốc của vật.
\item Tính tốc độ của vật sau $3\ s$ kể từ khi bắt đầu chịu lực.
\item Tính quãng đường $s$.
\item Nếu lực tác dụng trong $5\ s$ thay vì $3\ s$, quãng đường $s'$ vật đi được trong $5\ s$ là bao nhiêu?
\item Nếu khối lượng vật tăng gấp đôi, mà lực vẫn giữ nguyên, gia tốc mới bằng bao nhiêu?
\item Nếu muốn vật vẫn đạt tốc độ $6\ m/s$ sau $3\ s$ nhưng khối lượng tăng gấp đôi, phải thay đổi lực tác dụng như thế nào?
\end{listEX}
\loigiai{
\begin{itemize}
\item[1.] Gia tốc: $a = \dfrac{F}{m} = \dfrac{6}{2} = 3\ m/s^2$
\item[2.] Tốc độ sau $3\ s$: $v = at = 3 \cdot 3 = 9\ m/s$
\item[3.] Quãng đường: $s = \dfrac{1}{2} a t^2 = \dfrac{1}{2} \cdot 3 \cdot 9 = 13{,}5\ m$
\item[4.] Với $t = 5\ s$: $s = \dfrac{1}{2} \cdot 3 \cdot 25 = 37{,}5\ m$
\item[5.] Khối lượng tăng gấp đôi $\Rightarrow m' = 4\ kg \Rightarrow a' = \dfrac{6}{4} = 1{,}5\ m/s^2$
\item[6.] Để vật có $v = 6\ m/s$ sau $t = 3\ s$, cần $a = \dfrac{v}{t} = \dfrac{6}{3} = 2\ m/s^2$.

Mà $m = 4\ kg \Rightarrow F = ma = 4 \cdot 2 = 8\ N$
\end{itemize}
}
\end{vidu}
%\loai{Vận dụng trực tiếp công thức $F=ma$}
\begin{vidu}%[1P2-3.2Y][Loai1] 
Một quả bóng đang nằm yên trên mặt đất thì bị một cầu thủ đá bằng một lực $200\ N$. Bỏ qua mọi ma sát, gia tốc mà quả bóng thu được là $400\ m/s^2$. Khối lượng của quả bóng là
\choice
{\True $0,5\ kg$}
{$2\ kg$}
{$5\ kg$}
{$0,2\ kg$}
\haidong{3}
\loigiai{
Khối lượng của quả bóng là $m = \dfrac{F}{a} = \dfrac{200}{400} = 0,5\ kg$.
}
\haidong{8} \end{vidu}
%\loai{Tìm các thông số chuyển động}
\begin{vidu}%[1P2-3.2B][Loai1]
Một hợp lực $2\ N$ tác dụng vào một vật có khối lượng $2\ kg$ lúc đầu đứng yên, trong khoảng thời gian $2\ s$. Đoạn đường mà vật đó đi được trong khoảng thời gian đó là
\choice
{$1\ m$}
{$4\ m$}
{\True $2\ m$}
{$8\ m$}
\haidong{5}
\loigiai{


\begin{itemize}
\item Gia tốc mà vật thu được là $a = \dfrac{F}{m} = \dfrac{2}{2} = 1\ m/s^2$
\item Đoạn đường mà vật đó đi được trong khoảng thời gian đó là $s = {v_0}t + \dfrac{1}{2}at^2 = 0 + \dfrac{1}{2}.1.2^2 = 2\ m$.
\end{itemize}
}
\haidong{8} \end{vidu}
%\loai{Tính F}
\begin{vidu}%[1P2-3.2K][Loai1]
Một vật có khối lượng $2\ kg$ chuyển động thẳng nhanh dần đều từ trạng thái nghỉ. Vật đó đi được $200\ cm$ trong thời gian $2\ s$. Độ lớn hợp lực tác dụng của nó là
\choice
{$4\ N$}
{\True $2\ N$}
{$1\ N$}
{$100\ N$}
\haidong{5}
\loigiai{
\begin{itemize}
\item Gia tốc chuyển động của vật là $a = \dfrac{2s}{t^2} = \dfrac{2.2}{2^2} = 1\ m/s^2$
\item Độ lớn hợp lực tác dụng lên vật là $F = ma = 2.1 = 2\ N$.
\end{itemize}
}

\haidong{8} \end{vidu}
\begin{vidu}
Một lực không đổi tác dụng vào một vật có khối lượng $2,5\ kg$ làm vận tốc của nó tăng dần từ $2\ m/s$ đến $6\ m/s$ trong $2\ s$. Lực tác dụng vào vật có độ lớn bằng bao nhiêu niutơn?
\shortans[oly]{5}
\haidong{5}
\loigiai{\begin{itemize}
\item Gia tốc:
$
a=\dfrac{v_2-v_1}{t}=2\ m/s^2
$
\item Lực tác dụng: ${F}={ma}=5\ N$
\end{itemize}

} \haidong{8} \end{vidu}

%\loai{Thời gian va chạm}
\begin{vidu} %[1P2-3.2K][Loai1]
Một quả bóng, khối lượng $0,5\ kg$ đang nằm yên trên mặt đất. Một cầu thủ đá bóng với một lực $250\ N$. Thời gian chân tác dụng vào bóng là $0,02\ s$. Quả bóng bay đi với tốc độ bằng
\choice
{$0,01\ m/s$}
{$0,1\ m/s$}
{$2,5\ m/s$}
{\True $10\ m/s$}
\haidong{5}
\loigiai{      
\begin{itemize}
\item Ta có: $F = ma \Rightarrow a = \dfrac{F}{m} = \dfrac{250}{0,5} = 500\ m/s^2$.
\item Mà $a = \dfrac{v}{t} \Rightarrow v = at = 500.0,02 = 10$ m/s.
\end{itemize}
}

\haidong{8} \end{vidu}
%\loai{Lập tỉ lệ}
\begin{vidu}%[1P2-3.2K][Loai1]
Một xe tải chở hàng có tổng khối lượng xe và hàng là $4$ tấn, khởi hành với gia tốc $0,3\ m/s^2$. Khi không chở hàng xe tải khởi hành với gia tốc $0,6\ m/s^2$. Biết rằng lực tác dụng vào ô tô trong hai trường hợp đều bằng nhau. Khối lượng của xe lúc không chở hàng là
\choice
{$1,5$ tấn}
{$1,0$ tấn}
{\True $2,0$ tấn}
{$2,5$ tấn}
\haidong{5}
\loigiai{
Lực tác dụng trong hai trường hợp như nhau $\Rightarrow  F=m_1a_1=m_2a_2 \Rightarrow  4.0,3 = m_2.0,6 \Rightarrow  m_2 = 2$ tấn.
}
\haidong{8} \end{vidu}
%\loai{Hai lực tác dụng vào vật}
\begin{vidu}%[1P2-3.2K][Loai1] 
Một vật có khối lượng $2\ kg$, lúc đầu đang đứng yên. Nó bắt đầu chịu tác dụng đồng thời của hai lực $F_1=4\ N$  và $F_2=3\ N$. Biết góc giữa chúng là $30^\circ$. Quãng đường vật đi được sau $1,2\ s$ là
\choice
{$2,75\ m$}
{$2,65\ m$}
{$2,55\ m$}
{\True $2,45\ m$}
\haidong{5}
\loigiai{
\begin{itemize}
\item  Lực tổng hợp tác dụng lên vật: $F=\sqrt{F_1^2+F_2^2+2F_1F_2\cos 30^\circ}=6,8\ N$.
\item Gia tốc của vật: $a=\dfrac{F}{m}=\dfrac{6,8}{2}=3,4\ m/s^2 $.
\item Quãng đường vật đi được trong thời gian $t = 1,2$ s: $s=\dfrac{1}{2}a{t}^{2}=2,45\ m$.
\end{itemize}
}

\haidong{8} \end{vidu}
%\loai{Lực xiên góc}

%\loai{Đồ thị}
\begin{vidu}
\immini[thm]{Cho đồ thị biểu diễn mối liên hệ giữa các lực tác dụng lên một vật và gia tốc gây ra tương ứng như hình vẽ. Khối lượng của vật là
\choice
{$1,0\ kg$}
{$2,0\ kg$}
{\True $0,5\ kg$}
{$1,5\ kg$}}{\begin{tikzpicture}[draw=Mapcolor, very thick,scale=1,>=stealth']
\tkzDefPoints{0/0/o, 4.5/0/t, 0/2.5/d}
\tkzDrawSegments[very thick,->](o,t o,d)
\draw  [help  lines, xstep=0.5cm,ystep=0.5cm]  (0,0)  grid  (4,2);
\tkzLabelPoint[left](o){$O$}
\node[right] at (d){F (N)};
\node[above] at (t){a ($m/s^2$)};
\foreach  \x  [evaluate=\x as \z using \x] in  {1,2,...,4}
\node[below] at(\x,0){$\z$};
\foreach  \x  [evaluate=\x as \z using int(\x)] in  {1,2,...,4}
\draw(\x,-0.05)--(\x,0.05);
\foreach  \y  [evaluate=\y as \w using \y] in  {0.5,1.0,1.5,2.0}
\node[left] at(0,\y){$\w$};
\foreach  \y  [evaluate=\y as \z using int(\y)] in  {0.5,1.0,...,2.0}
\draw(-0.05,\y)--(0.05,\y);
\draw[Mapcolor,line width=1.2pt](0,0)--++(26.7:4);
\draw[fill=white](1,0.5)circle(2pt);
\draw[fill=white](2,1)circle(2pt);
\draw[fill=white](3,1.5)circle(2pt);
\end{tikzpicture}}
\haidong{5}
\loigiai{Khối lượng của vật: $m=\dfrac{F}{a}=\dfrac{1}{2}=0,5\ kg$}
\haidong{8} \end{vidu}	
%\loai{Tổng hợp}
\begin{vidu}%[1P2G3-2][Loai1]
Xe có khối lượng $800\ kg$ đang chuyển động thẳng đều thì hãm phanh chuyển động chậm dần đều. Biết quãng đường đi được trong giây cuối cùng của chuyển động là $1,5\ m$. Lực hãm của xe có độ lớn là
\choice
{\True $2400\ N$}
{$260\ N$}
{$240\ N$}
{$2600\ N$}
\haidong{5}
\loigiai{
\begin{itemize}
\item Giả sử xe đang chuyển động với vận tốc $v_0$. 
\item Sau khoảng thời gian t thì xe dừng lại.
\item Quãng đường xe đi được trong giây cuối cùng là: 
\begin{align*}
\Delta s = {v_0}t + \dfrac{1}{2}at^2 - {v_0}\left( {t - 1} \right) - \dfrac{1}{2}a{\left( {t - 1} \right)^2} = 1,5 \Rightarrow {v_0} + at - \dfrac{a}{2} = 1,5.
\end{align*}
\item Do $v_0 + at = 0\Rightarrow \dfrac{-a}{2} = 1,5\Rightarrow a = -3\ m/s^2$.
\item Lực hãm của xe là $F = ma = 800.(-3) = -2400$ N.
\end{itemize}
}	
\haidong{8} \end{vidu}
%\loai{Thay đổi khối lượng}
\begin{vidu}
Một lực ${F}$ truyền cho vật có khối lượng ${m}_1$ một gia tốc bằng $6\ m/s^2$, truyền cho vật khác có khối lượng ${m}_2$ một gia tốc bằng $3\ m/s^2$. Nếu đem ghép hai vật đó lại thành một vật thì lực đó truyền cho vật ghép một gia tốc bằng bao nhiêu $m/s^2$?
\shortans[oly]{2}
\haidong{5}\loigiai{\begin{itemize}
\item 	$
\text{Ta có:} {a}_1=\dfrac{F}{m_1} ;\ {a}_2=\dfrac{F}{m_2} \Rightarrow {m}_1=\dfrac{F}{a_1} ;\ {m}_2=\dfrac{F}{a_2} ; \\
\Rightarrow {a}=\dfrac{F}{m_1+m_2}=\dfrac{F}{\dfrac{F}{a_1}+\dfrac{F}{a_2}}=\dfrac{a_1 a_2}{a_1+a_2}=2\ m/s^2 .
$
\end{itemize}

} \haidong{8} \end{vidu}
\begin{vidu}%[1P2-3.2K][Loai1]
\immini[thm]{Vật khối lượng $m = 2\ kg$ đặt trên mặt sàn nằm ngang và được kéo nhờ lực như hình, $\vec F$  hợp với mặt sàn góc $\alpha  =  60^\circ$  và có độ lớn $F = 2\ N$. Bỏ qua ma sát. Độ lớn gia tốc của m khi chuyển động là  
\choice
{$1\ m/s^2$}
{$0,85\ m/s^2$}
{\True $0,5\ m/s^2$}
{$0,45\ m/s^2$}}{\begin{tikzpicture}[scale=1,thick,>=stealth',draw=Mapcolor]
\node  at  (0,0.32)  [shape=rectangle,draw=red,pattern=dots,minimum height=0.6cm,minimum width=1.4cm, rounded corners=1pt]  {};
\draw(-1.5,0)--(1.5,0);
\draw[very thick,Mapcolor,->](0.7,0.3)--++(60:2)node[right,black]{$\vec{F}$}coordinate(f);
\draw[fill=white](0.7,0.3)circle(1.2pt);
\draw[dashed](0.7,0.3)coordinate(o)--(1.7,0.3)coordinate(fx);
\tkzMarkAngles[size=0.4cm](fx,o,f)
\tkzLabelAngles[pos=0.8](fx,o,f){$60^\circ$}
\fill[pattern=north east lines](-1.5,0)--(1.5,0)--++(-90:0.1)--++(180:3);
\end{tikzpicture} }
\haidong{5}
\loigiai{
\immini{\begin{itemize}
\item Ta phân tích lực như hình vẽ.
\item Thành phần lực $F_x$ có tác dụng làm cho vật chuyển động trên mặt sàn
\begin{align*}
F_x = F\cos {60^\circ} = 2.\cos {60^\circ} = 1\ N
\end{align*}
\item Độ lớn gia tốc của m khi chuyển động là $a = \dfrac{F}{m} = \dfrac{1}{2} = 0,5\  m/s^2$.
\end{itemize}}{\begin{tikzpicture}[scale=1,thick,>=stealth']
\tkzDefPoints{0/1/o}
\node  at  (0,0.32)  [shape=rectangle,draw=red,pattern=dots,minimum height=0.6cm,minimum width=1.4cm, rounded corners=1pt]  {};
\draw(-1.5,0)--(1.5,0);
\draw[very thick,blue,->](0.7,0.3)--++(60:2)node[right,black]{$\vec{F}$}coordinate(f);
\draw[fill=white](0.7,0.3)circle(1.2pt);
\draw[very thick,blue,->](0.7,0.3)--++(0:1)node[below right,black]{$\vec{F}_y$}coordinate(fx);
\draw[very thick,blue,->](0.7,0.3)coordinate(o)--++(90:1.73)node[above,black]{$\vec{F}_x$}coordinate(fy);
\draw[dashed](fx)--(f)--(fy);
\tkzMarkAngles[size=0.4cm](fx,o,f)
\tkzLabelAngles[pos=0.8](fx,o,f){$60^\circ$}
\fill[pattern=north east lines](-1.5,0)--(1.5,0)--++(-90:0.1)--++(180:3);
\end{tikzpicture} }
}

\haidong{8} \end{vidu}

%\end{comment}